%styl klasy z Polskimi Normami oprac. Marcin Wolinski
\documentclass[a4paper,titleauthor]{mwart} 

\usepackage{polski}
\usepackage[utf8]{inputenc}
\usepackage{graphicx} %pakiet do wstawiania grafiki
\usepackage[hyphens]{url} %pakiet do wstawiania linkow
%\usepackage[hidelinks,breaklinks]{hyperref}
\usepackage{authblk}%pakiet do tworzenia afiliacji
\usepackage{tabularx}%pakiet do tabel
\usepackage[a4paper, left=2cm, right=2cm, top=3cm, bottom=3cm]{geometry}
\usepackage{listings}
\usepackage{placeins}%pakiet do kontroli umieszczania obiektow
\usepackage{hyperref}%pakiet do m.in. kolorowania linkow
\usepackage[]{xcolor}
\usepackage{float}
\usepackage{minted}


\usepackage[tablegrid,owncaptions]{vhistory}
\renewcommand{\vhhistoryname}{Historia zmian}
\renewcommand{\vhversionname}{Wersja}
\renewcommand{\vhdatename}   {Data}
\renewcommand{\vhauthorname} {Autor}
\renewcommand{\vhchangename} {Opis zmian}

\renewcommand\figurename{Rys.}%skrocony podpis
\renewcommand\lstlistingname{Wydruk}


%------------------------------------------------------------------------
% Dane do strony tytułowej

\title{{\Huge  Projekt BDBT}\\ - \\{\Large Zespół A2}\\ }

\author{Szymon Boniuk XXX \and Jakub Celiński 337318}

\affil{Politechnika Warszawska, Wydział EiTI}
\date{\today}

%------------------------------------------------------------------------
% Początek dokumentu
\begin{document}
	
	%Automatycznie generowany tytuł dokumentu
	\maketitle
	
	\newpage
	\tableofcontents
	\newpage
	%------------------------------------------------------------------------
	\section{Zakres i cel projektu (opis założeń funkcjonalnych projektowanej bazy danych)}
	
	\section{Definicja systemu}
	\subsection{Perspektywy użytkowników}
	
	\section{Model konceptualny}
	\subsection{Definicja zbiorów encji określonych w projekcie (decyzje projektowe)}
	\subsection{Ustalenie związków między encjami i ich typów}
	\subsection{Określenie atrybutów i ich dziedzin}
	\subsection{Dodatkowe reguły integralnościowe (reguły biznesowe)}
	\subsection{Klucze kandydujące i główne (decyzje projektowe)}
	\subsection{Schemat ER na poziomie konceptualnym}
	\subsection{Problem pułapek szczelinowych i wachlarzowych – analiza i przykłady}
	
	\section{Model logiczny}
	\subsection{Charakterystyka modelu relacyjnego}
	\subsection{Usunięcie właściwości niekompatybilnych z modelem relacyjnym - przykłady}
	\subsection{Proces normalizacji – analiza i przykłady}
	\subsection{Schemat ER na poziomie modelu logicznego}
	\subsection{Więzy integralności}
	\subsection{Proces denominacji – analiza i przykłady}
	
	\section{Faza fizyczna}
	\subsection{Projekt transakcji i weryfikacja ich wykonalności}
	\subsection{Strojenie bazy danych – dobór indeksów}
	\subsection{Skrypt SQL zakładający bazę danych}
	\subsection{Przykłady zapytań i poleceń SQL odnoszących się do bazy danych}
	\newpage

\bibliographystyle{unsrt} % plplain plabbrv plalpha
\bibliography{BDBT}

\end{document}

%\begin{figure}[H]
%	\centering
%	\includegraphics[scale = 0.6]{nazwa.png}
%	\caption{PODPIS}
%	\label{fig:PODPIS}
%\end{figure}